\section{Parametric-Source Transformations}
\label{sec:param-source}

Two transformations take a Parametric carrier as their source: Evaluator-Driven Optimization moves the judgment capability fixed in evaluator parameters into Policy parameters, and Parametric Externalization externalizes the decision or evaluation capability already internalized in the weights into Tokenized artifacts that other systems can consume. The source of both pathways is the product of another pathway, not direct supervision from a raw interaction log.

\subsection{Evaluator-Driven Optimization: Evaluator-to-Policy Transformation}
\label{sec:evaluator-driven}

Evaluator-Driven Optimization transfers the evaluative capability a parametric evaluator has learned into Policy parameters through training: the Policy as actor produces behavior, the evaluator judges that behavior, and the resulting judgment serves as the training signal that updates the Policy weights. The source of this pathway is the Evaluator parameter rather than the judged $(c,a,o)$---the judged behavior is only the object through which the evaluator expresses its judgment, and the experience truly transferred was already compressed into the evaluator weights back in the Evaluator Internalization stage; the Evaluator parameter is both the product of that pathway and the starting point of this one. The boundary with Policy Internalization lies in the source of the learning signal: the latter writes raw trajectories, human demonstrations, or environment success signals into the weights directly, whereas the former takes the evaluator's judgment of the Policy's behavior as the signal. The criterion requires that the Policy weights are indeed updated by the evaluator signal; work in which the evaluator only modulates decoding, reranking, best-of-N, tree search, or MCTS without updating the weights does not belong to this pathway.

\subsubsection{Discriminative Outcome Evaluator-to-Policy Transfer}

A discriminative outcome evaluator gives a comparable scalar score, or a chosen/rejected preference, over a complete output or a complete trajectory, and the policy writes these signals into its parameters. By how the policy absorbs the signal, we distinguish three classes: Reward Maximization treats the scalar as a reward and maximizes it through RL, Preference Contrast lets the preference drive contrastive optimization directly, and Reward-Ranked Selection uses the score to filter samples for supervised writing.

\paragraph{Reward Maximization.} The evaluator's scalar score acts as a reward, which the policy maximizes through an RL algorithm such as PPO, GRPO, or REINFORCE. The RL-CAI stage of Constitutional AI~\cite{Bai22b} has a constitution-guided AI judge give preferences over complete answers, trains a preference model, and then has the policy maximize its score through RLAIF. RLAIF~\cite{Lee23b} gives two uses of AI feedback: canonical RLAIF first trains an RM from AI preferences, and d-RLAIF takes the judge's score directly as an online reward. A further work removes the ground-truth label entirely: an LLM judge answers a meta-question whose YES probability is synthesized into a scalar reward maximized through GRPO/CISPO~\cite{Ren26b}. Two works target the reward hacking that a policy invites by optimizing an imperfect evaluator over the long run. The first~\cite{Ack26} adds gradient regularization and a reference reset to a Dr.GRPO-style objective, suppressing the policy's exploitation of sharp peaks in the proxy reward; Policy Filtration~\cite{She24c} admits only the RM's most reliable high- and low-scoring samples into the PPO buffer, easing overfitting to a noisy RM. The evaluators of FAPO~\cite{Din25b}, a related method~\cite{Che25v}, and OS-Themis~\cite{Li26p} have step-level diagnosis but compress it into a trajectory scalar before entering RL, landing on the outcome side (bordering Section~\ref{sec:disc-process-transfer}): FAPO folds the localization of the first invalid step into a penalty on the whole rollout, the related method~\cite{Che25v} computes a partial-correctness scalar for a trajectory from first-error localization to activate the gradient of an all-negative group, and OS-Themis takes a trajectory-level reward as its main online-RL signal.

\paragraph{Preference Contrast.} The evaluator's chosen/rejected preference is not distilled into a separate reward but enters a contrastive objective such as DPO, IPO, or SLiC to update the policy directly. OAIF~\cite{Guo24} is an on-policy instance: an online LLM annotator gives preferences over response pairs sampled from the current policy, fed straight into DPO/IPO/SLiC, so the supervision distribution always tracks the policy and the RM never drifts off-distribution as training proceeds. Self-Rewarding Language Models~\cite{Yua24} and Meta-Rewarding Language Models~\cite{Wu24d} have the evaluator and policy share one set of parameters: the former has the same model act as both generator and judge, building chosen/rejected pairs from self-scoring for iterative DPO; the latter adds a meta-judge layer that improves the judge's own scoring quality, forming an actor/judge/meta-judge co-evolution. Two works improve the robustness of preference optimization: one uses a Bayesian router to select dynamically among several RMs by query, feeding the routed preference pairs into online DPO~\cite{Wu25x}; the other, in an offline setting, distills the RM's reward difference into the policy's implicit reward through L2 distillation and a forward-KL regularizer, improving stability under label bias and distribution shift~\cite{Fis24}.

\paragraph{Reward-Ranked Selection.} The evaluator score is used only to select high-scoring samples, which are then written into the policy by SFT, approaching the alignment effect of RL at lower compute. RAFT~\cite{Don23} has an RM score several answers to the same prompt and keeps the best-of-K for SFT. BOND~\cite{Ses24} distills the entire best-of-N distribution---not just the high-scoring samples---into a single-sample policy through forward and backward KL. BoNBoN~\cite{Gui24} characterizes the best-of-n distribution offline and admits both the best and worst samples through SFT-BoN and IPO-BoN. The critic LLM of \cite{Gon24b} scores a whole trajectory and takes the top $p\%$ as pseudo-demonstrations for iterative SFT alongside general data.

Mixture of Judges~\cite{Xu24d} uses a mixture of signals from a calibrated RM, an LLM judge, and a rule-constrained judge to instantiate CRPG, CODPO, and CRRAFT respectively, spanning all three classes above; its value for multi-evaluator de-biasing is discussed in Section~\ref{sec:edo-discussion}.

\subsubsection{Generative Outcome Evaluator-to-Policy Transfer}

A generative outcome evaluator produces natural-language content---critique, failure explanation, or revised response---over a complete output or trajectory, and takes that content rather than a scalar as the policy's main supervision.

The SL-CAI stage of Constitutional AI~\cite{Bai22b} is an early instance: the feedback model generates a critique and a revision over a complete answer by the constitution, and the policy SFTs on the revision directly. ECHO~\cite{Li26l} is a more recent full instance: a linguistic critic updated in step with the actor generates a score-aware critique over a complete rollout, the actor produces a refined rollout accordingly, the policy optimizes those refined trajectories with GRPO conditioned on the critique, and the critic co-evolves through a dual-track GRPO on the same batch. The value of generative feedback naturally falls to the step level---once an evaluator is willing to generate fine-grained explanations, using them at process granularity makes better use of the information than at outcome granularity---so generative outcome evaluators mostly migrate toward Section~\ref{sec:gen-process-transfer}.

\subsubsection{Discriminative Process Evaluator-to-Policy Transfer}
\label{sec:disc-process-transfer}

A discriminative process evaluator gives a comparable step-level signal over an intermediate step, a local action, a reasoning prefix, or a state-action pair, and the policy updates its parameters accordingly. The core motivation is credit assignment for a long-horizon agent: from a terminal reward alone, the policy struggles to tell which local behaviors to reinforce or avoid. By how the step-level signal is produced, we distinguish three classes: Step-Value Critic trains a critic to estimate a step-level value, Step-Quality Scoring has the evaluator score each step's quality directly, and Reward Redistribution spreads an existing trajectory-level reward across the steps.

\paragraph{Step-Value Critic.} This class trains a dedicated critic to estimate the value or advantage of each step as a baseline for credit assignment. SWEET-RL~\cite{Zho25p} trains a turn-level critic in multi-turn collaborative reasoning, learning the advantage of each turn's action directly from training-time privileged information, then distilling the turn-wise credit assignment into the actor through DPO. The Agentic-Q model of \cite{Wan26s} estimates the future return of a state-action in a GUI setting and feeds the step-wise value into an SWPO-style update.

\paragraph{Step-Quality Scoring.} The evaluator scores each step, or each dimension, directly, and the policy writes the signal in through RL reward, filtering, or preference. PPRM~\cite{Xu25h} targets non-verifiable multi-turn search, scoring each turn's correctness, relevance, and coherence with a principle-based PRM and calibrating through ReNorm and a terminal ORM before feeding into PPO+GAE. Language Feedback Model~\cite{Zho24e} and \cite{Pan24b} both turn per-step quality judgments into a filtering signal: the former judges whether each action is productive and does behavior cloning on the state-actions marked desirable; the latter does Filtered BC on the results of per-step progress filtering. AdaRubric~\cite{Din26e} and \cite{Luo25f} both build the quality score into a preference: the former generates a rubric for the task and then outputs a 1--5 score for each step and dimension, aggregating into trajectory-level DPO and step-level PPO; the latter has a judge score single-step GUI actions for computer use on a 0--100 scale and builds step-level preference pairs from high and low scores to train UI-TARS-2B.

\paragraph{Reward Redistribution.} This class builds no extra step-level labels but spreads an existing trajectory-level reward across tokens or rounds. One method takes the token attribution of a standard RM's last-layer attention and redistributes the terminal scalar reward into a token-level dense reward for PPO~\cite{Cha24b}. CW-GRPO~\cite{Wan26ac} has a rubric-based judge give two binary signals---retrieval utility and reasoning correctness---for each search round, synthesizing a round-level contribution weight that redistributes the trajectory-level advantage by round.

\subsubsection{Generative Process Evaluator-to-Policy Transfer}
\label{sec:gen-process-transfer}

A generative process evaluator produces explanatory, diagnostic, or corrective content over intermediate steps, and turns that content into the policy. Unlike a discriminative process evaluator, which answers ``is this step good,'' a generative one answers ``why is this step bad and how should it change.'' By the form of the generated content, we distinguish two classes: Step Critique-and-Refine critiques each step and produces a refined step, and Failure-Driven Revision localizes errors or gives instructions and then internalizes the correction.

\paragraph{Step Critique-and-Refine.} The evaluator critiques each step, the actor produces a refined step accordingly, and the refined content is written into the policy as the alignment target. SRR-Judge~\cite{Zha26am} outputs an explanation, a 1--5 rating, and a refined thought/action for each thought-action step of a search agent, with the rating handling selection and the refined step serving as the generative alignment target, written into the policy through iterative RFT. Natural Language Actor-Critic~\cite{Hon25c} uses a generative language critic to produce, for each state-action, a natural-language critique of its merits and future direction; the actor generates a refined action conditioned on the critique, which is then distilled back into the base policy.

\paragraph{Failure-Driven Revision.} The evaluator localizes an error or gives a corrective instruction, and the corrected trajectory or action is written into the policy through internalization. Agent-R~\cite{Yua25c} has the actor itself serve as verifier, localizing the transition point of a bad trajectory at a search-tree branch, inserting a revision signal between the erroneous prefix and the correct continuation, and iteratively SFT-ing the ``reflect-and-correct'' instances. OpenClaw-RL~\cite{Wan26ad} extracts directive textual hints from user utterances or environment signals and uses On-Policy Distillation to turn the log-prob gap between a hint-enhanced teacher and the student into a token-level directional update; its sequence-level binary RL branch falls outside this cell.

\subsubsection{Discussion}
\label{sec:edo-discussion}

What the works under Evaluator-Driven Optimization share is the insertion of a learned evaluator between the experience and the policy, and both the power and the limits of this pathway are set by that intermediary. It makes alignment scalable---once ``what output is better'' is compressed offline into the evaluator weights, the marginal cost of sampling, filtering, and contrasting is very low, which is the precondition for the RLHF toolchain to scale horizontally---and it extends the trainable range into the non-verifiable region a rule-based verifier cannot reach: open dialogue, long-horizon search, and multi-dimensional agentic behavior give no clean reward from the environment, and the evaluator is almost the only means of turning a hard-to-quantify preference into an optimizable signal. The same intermediary sets a ceiling: the policy cannot exceed the evaluator; optimizing a fixed, imperfect evaluator continually drives the policy toward what the evaluator deems good rather than the true goal, which makes reward hacking an intrinsic risk of this pathway, recurring across the scalar/critique and outcome/step forms.

The pathway unfolds along two axes, each carrying a tension. The granularity axis (outcome/process) sets the precision of credit assignment: an outcome signal is simple but its long-horizon credit is coarse; a process signal gives dense supervision but the reliability and consistency of step labels are harder to guarantee. The form axis (discriminative/generative) sets the trade-off between information density and verifiability: a discriminative signal is low in density but comparable and easy to filter and rank; a generative signal is high in density but hard to verify, and its boundary with tokenized carriers such as a refined trajectory is blurred, often pulling a method's membership toward Policy Internalization or Parametric Externalization. The two axes evolve in the same direction---from discriminative to generative, from outcome to process---raising information density while pushing the question of whether the evaluator itself is reliable ever further forward.

The reliability of the evaluator is the core constraint of this pathway and its largest open problem. Existing mitigations each have an emphasis---co-evolution lets the evaluator update with the policy~\cite{Yua24, Wu24d, Li26l}, ensembling and routing spread the single-point bias across several evaluators~\cite{Wu25x, Xu24d}, and regularization and filtration constrain the optimization process~\cite{Ack26, She24c}---but all can only suppress, not eliminate, the gap between the evaluator and the true goal. Verifiability is also asymmetric: a discriminative signal can still be checked indirectly through held-out preference consistency, whereas the correctness of a generative step-level critique remains dependent on humans or a stronger model to this day, and a wrong critique misleads more than a wrong scalar. Placed back among its neighbors, Policy Internalization's supervision is grounded directly but passes the noise straight in, whereas Evaluator-Driven Optimization trades a layer of grounding for a steadier, more comparable signal, the two faring differently when environment signals are abundant versus sparse; and the evaluator is itself the product of Evaluator Internalization, so the P4 $\rightarrow$ P6 chaining forms a complete evaluator-based pipeline, whose composite form and cross-pathway systematic comparison are left to Section~\ref{sec:composite} and Section~\ref{sec:synthesis} respectively.

\subsection{Parametric Externalization: Parametric-to-Tokenized Transformation}
\label{sec:parametric-externalization}

Parametric Externalization externalizes the experience already internalized in Policy or Evaluator parameters into Tokenized artifacts that other systems can use. The externalized products fall into two kinds, matching the roles of the two ends of parameters: the Policy end produces an agent trajectory (a complete record of decision behavior), and the Evaluator end produces an evaluative annotation (a judgment of behavior quality, including a correctness/quality label, a natural-language critique, or a preference annotation). Unlike the other pathways, which process a trajectory that already exists, this pathway has its artifact generated on the spot by the model that carries the parameters---the same parameters can produce countless trajectories under different prompts and environments, so the transformation operator samples from $p(e \mid \text{prompt}, \text{env})$ rather than re-encoding a given $e$.

By which end of parameters the externalized product comes from, this pathway splits into two: Demonstration Externalization (Section~\ref{sec:demo-ext}) externalizes the Policy's decision behavior, and Evaluative Supervision Externalization (Section~\ref{sec:eval-sup-ext}) externalizes the Evaluator's judgment capability.

\subsubsection{Demonstration Externalization}
\label{sec:demo-ext}

Demonstration Externalization externalizes the implicit decision capability of a parametric Policy or teacher agent into agent trajectories that a student can imitate, with downstream SFT, behavioral cloning, or imitation training a new model. The faithfulness of the externalized trajectory depends on the environment the teacher generates in: when the teacher interacts in a real environment, the observation is the environment's objective return and the trajectory is most strongly grounded; when the environment is simulated by a model, or omitted entirely, the observation is generated by the model itself, the trajectory comes closer to a pure parametric read-out, the risks of pseudo-success and hallucinated observation rise, and the dependence on generation-stage verification grows accordingly. We split the work into Real-Environment Synthesis, Simulated-Environment Synthesis, and Constructive Synthesis.

\paragraph{Real-Environment Synthesis.} The teacher executes a task in a real, interactive environment, the observation comes from the environment's objective return, and grounding is strongest. The representative in the web domain is Explorer~\cite{Pah25}: GPT-4o discovers tasks autonomously on real web pages, executes multi-step interactions, and---after multi-stage refinement and a verifier---produces a multimodal trajectory with page state, action sequence, and summary. Extending along this line, InSTA~\cite{Tra25} pushes to internet scale, AutoSurfer~\cite{Fai26} first identifies a site's functional space and then organizes it into task tuples, SynthAgent~\cite{Wan25ar} does element-level exploration before a global refine, and A3-SYNTH~\cite{Lu26} uses Agents-as-Annotators to produce successful trajectories with structured reasoning; a further work uses a judge to extract constraints and does prefix extraction and hindsight relabeling by satisfaction rate, making partially successful trajectories trainable too~\cite{Log26b}. In the computer and mobile domains, Fara-7B~\cite{Awa25} produces screenshot-thought-action trajectories through FaraGen and filters them with a triple verifier, OpenMobile~\cite{Che26d} has an expert take over and correct learner deviations and keeps the corrected trajectory and step-level CoT, and AgentSynth~\cite{Xie25e} first generates and verifies simple subtasks in environments such as OSWorld, then composes longer compound tasks while keeping the full action trajectory and visual evidence. In the tool and MCP domains, TOUCAN~\cite{Xu25o} synthesizes tasks and executes full tool-use rollouts in nearly 500 real MCP environments, forming a 1.5M-scale corpus through task-level and trajectory-level double filtering. ANCHOR~\cite{Wei26b} finds branch points along a few high-quality seed trajectories, poses new tasks at the branching state and continues, forming tree-shaped GUI data with a shared prefix and branching tails.

\paragraph{Simulated-Environment Synthesis.} The environment itself is played by a model, and the observation is generated by a simulator. This group is most concentrated in the tool-call domain: ToolAlpaca~\cite{Tan23} has ChatGPT play user, assistant, and tool executor to simulate tool trajectories, ToolMind~\cite{Yan25ab} has three kinds of agent simulate multi-turn trajectories with reasoning and does fine-grained filtering, ToolACE~\cite{Liu24o} first expands the API pool and then synthesizes tool-call dialogues of varied complexity, and APIGen-MT~\cite{Pra25b} fixes a task blueprint and ground-truth action first and then synthesizes multi-turn function-calling dialogue through a simulated agent-human interplay. In the mobile GUI domain, Learning with Challenges~\cite{Kan26b} first profiles the student's capability and then uses explorer/supervisor/synthesizer VLM roles to generate difficulty-adaptive goal-conditioned trajectories in an emulator. One group makes the ``environment itself'' the main product: \cite{Wan25as} uses GPT-4o-mini as both world simulator and teacher, producing both trajectories and an extensible UI state-transition simulator; Mock Worlds~\cite{Lu26j} synthesizes a workflow, a virtual tool ecosystem, and a rubric from a persona; ASTRA~\cite{Tia26} produces SFT trajectories on one track and compiles tool implementations into a sandbox-verifiable RL arena on another; and a further work self-evolves on a few seeds and then perturbs queries and tools to build an RL-ready mock environment~\cite{Xu26e}.

\paragraph{Constructive Synthesis.} This end has almost no interactive environment; the trajectory is constructed in reverse around a pre-given answer, evidence, or problem structure, coming closest to a pure parametric read-out. EviPath~\cite{Li25at} back-derives an abductive reasoning path from gold evidence and then generates a RAG trajectory with planning and evidence-grounded answering; RAGShaper~\cite{Tao26} builds an information tree with distractors and multi-hop questions and induces an agentic RAG trajectory from the teacher in a controlled retrieval environment; and APIGen~\cite{Liu24n} generates queries and function-calling answers from real executable APIs and filters them through three-stage checking (it keeps an execution check, but its interaction is shallow and it leans toward generation overall).

\subsubsection{Evaluative Supervision Externalization}
\label{sec:eval-sup-ext}

Evaluative Supervision Externalization externalizes the implicit evaluation capability of a parametric Evaluator (judge, verifier, critic, reward model) into Tokenized evaluative annotation, used to train a student, distill a lightweight evaluator, or filter data. By form, the products fall into three kinds: a correctness/quality label at the step or outcome level, a natural-language critique and failure diagnosis, and a preference annotation.

One work first has a strong computer-use agent produce noisy trajectories, then has o4-mini grade each step for correctness and optimality with screenshot, action, and alternative analysis~\cite{He25g}; the results are materialized into two independent datasets, WebSTAR and WebSCORE, the former training a student agent and the latter distilling a lightweight StepRM. Step-GUI~\cite{Yan25x} first establishes a quality anchor with a trajectory-level verifier, then has a strong multimodal model transcribe the verification results into seven kinds of structured step-level supervision---progress tracking, state summary, effect prediction, and others---used together with the original trajectory for cold-start fine-tuning and RLVR.

\subsubsection{Discussion}

The value of Parametric Externalization lies at two levels. The first is data production: it replaces item-by-item human annotation with a teacher's automatic generation, lowering unit cost and synthesizing the needed experience directionally along dimensions such as curriculum, difficulty, domain coverage, and failure pattern, so data scale and coverage are no longer bound by human throughput. The second is capability transfer: the implicit capability of an expensive large teacher can be externalized and then distilled into a set of lightweight, domain-specialized students, trading a recurring high inference cost for a one-time synthesis cost.

The cost mostly derives from one fundamental constraint: on the task-capability dimension, the student inherits the teacher's ceiling, and what this pathway buys is cost, inference efficiency, and domain specialization, not the capability frontier itself. Teacher hallucination, pseudo-successful trajectories, and non-executable steps are absorbed by the student as correct supervision, a risk that rises along the environment-synthesis axis---at the real-environment end the verifier need only check task success, while at the simulated and constructive end it must also check the credibility of the observation itself; distribution narrowing further magnifies a single teacher's style and preference, which multi-teacher mixing can ease but not remove. The verifier is therefore both the backbone of this pathway and its weakest point: a false positive lets a pseudo-success slip into the student and a false negative kills a high-quality sample, and scaling an evaluation-grade verifier in step with generation scale is the open engineering problem of this pathway.

On the experience-source dimension, Parametric Externalization is the main route by which \{teacher\}-sourced experience enters the training corpus---the other pathways mostly process experience produced by the agent itself or by humans, whereas the implicit capability of a teacher model must be externalized through this pathway before it can become explicit experience usable by other systems.
